% Part: first-order-logic
% Chapter: proof-systems
% Section: introduction

\documentclass[../../../include/open-logic-section]{subfiles}

\begin{document}

\iftag{FOL}
      {\olfileid{fol}{prf}{int}}
      {\olfileid{pl}{prf}{int}}

\olsection{प्रस्तावना}

तर्कप्रणालींमध्ये सामान्यतः चिन्हार्थमीमांसा आणि !!a{derivation}
पद्धती हे दोन्ही घटक असतात. चिन्हार्थमीमांसेत सत्यता, पूर्ततायोग्यता,
वैधता आणि तार्किक निष्पन्नता अशा संकल्पनांचा विचार केला जातो.
तार्किक निष्पन्नता आणि वैधता प्रस्थापित करण्यासाठी निव्वळ विन्यासात्मक
रीत उपलब्ध करून देणे हा !!{derivation} पद्धतींचा उद्देश असतो. त्या
निव्वळ विन्यासात्मक असतात याचा अर्थ असा की अशा पद्धतीतील
!!{derivation} ही एक सांत विन्यासात्मक वस्तू असते; सहसा ती
!!{sentence}s किंवा !!{formula}s यांचा अनुक्रम (किंवा इतर एखादी सांत
मांडणी) असते. !!{sentence}s किंवा !!{formula}s यांचा कोणताही दिलेला
अनुक्रम किंवा मांडणी ``योग्य'' आहे, हे यांत्रिक रीतीने पडताळता येते,
हा चांगल्या !!{derivation} पद्धतींचा गुणधर्म असतो.

प्रथम-क्रम तर्कशास्त्रासाठीच्या सर्वांत सोप्या (आणि इतिहासात प्रथम
आलेल्या) !!{derivation} पद्धती \emph{स्वयंसिद्धकीय} होत्या. अशा
पद्धतीत !!{formula}s चा एखादा अनुक्रम तेव्हाच !!a{derivation} मानला
जातो, जेव्हा त्यातील प्रत्येक स्वतंत्र !!{formula} एकतर
``स्वयंसिद्धकांच्या'' एका निश्चित संचात असते किंवा त्या अनुक्रमात
त्याच्या आधी आलेल्या !!{formula}s पासून ठरावीक संख्येतील ``निगमन
नियमां''पैकी एखाद्या नियमाने निष्पन्न होते. शिवाय, !!a{formula}
स्वयंसिद्धक आहे का आणि ते इतर !!{formula}s पासून एखाद्या निगमन
नियमानुसार योग्य रीतीने निष्पन्न होते का, हे यांत्रिक रीतीने पडताळता
येते. स्वयंसिद्धकीय !!{derivation} पद्धतींचे वर्णन करणे सोपे असते आणि
अधिउपपत्तीच्या पातळीवरही त्या हाताळणे सोपे असते; पण त्यांतील
!!{derivation}s वाचणे आणि समजून घेणे कठीण असते, तसेच त्या तयार करणेही
कठीण असते.

!!{derivation}s तयार करणे सोपे व्हावे किंवा पूर्ण झालेल्या
!!{derivation}s समजून घेणे सोपे व्हावे, या उद्देशाने इतर
!!{derivation} पद्धती विकसित करण्यात आल्या आहेत. नैसर्गिक निगमन,
सत्यता-वृक्ष---ज्यांना टॅब्लो सिद्धता असेही म्हणतात---आणि क्रमवर्ती कलन
ही त्यांची उदाहरणे आहेत. काही !!{derivation} पद्धती विशेषतः
यांत्रिकीकरण डोळ्यांसमोर ठेवून रचल्या जातात. उदाहरणार्थ, निराकरण पद्धत
संगणकीय प्रणालीत कार्यान्वित करणे सोपे असते (पण तिच्यातील
!!{derivation}s समजून घेणे जवळजवळ अशक्य असते). यांपैकी बहुतेक इतर
!!{derivation} पद्धती !!{derivation}s ना अनुक्रमांऐवजी !!{formula}s चे
वृक्ष म्हणून दर्शवतात. त्यामुळे !!a{derivation} चा कोणता भाग इतर
कोणत्या भागांवर अवलंबून आहे, हे पाहणे सोपे होते.

म्हणून प्रथम-क्रम तर्कशास्त्रासारखी एखादी तर्कप्रणाली घेतल्यास,
!!a{sentence} हे \emph{प्रमेय} असणे म्हणजे काय आणि !!a{sentence} हे
इतर काही वाक्यांपासून !!{derivable} असणे म्हणजे काय, यांची
वेगवेगळ्या !!{derivation} पद्धती वेगवेगळी नेमकी स्पष्टीकरणे देतील. हे
कसेही केले (स्वयंसिद्धकीय !!{derivation}s, नैसर्गिक निगमने, क्रमवर्ती
!!{derivation}s, सत्यता-वृक्ष किंवा निराकरण-खंडने यांच्या साहाय्याने),
तरी हे संबंध वैधता आणि तार्किक निष्पन्नता या चिन्हार्थविषयक संकल्पनांशी
जुळावेत, अशी आपली अपेक्षा असते. ``$!A$~हे प्रमेय आहे'' यासाठी
$\Proves !A$ आणि ``$!A$~हे~$\Gamma$ पासून !!{derivable} आहे'' यासाठी
``$\Gamma \Proves !A$'' लिहू या. $\Proves$~ची व्याख्या कोणतीही असो,
ती $\Entails$~शी जुळावी, म्हणजे:
\begin{enumerate}
\item $\Proves !A$ जर आणि फक्त जर $\Entails !A$
\item $\Gamma \Proves !A$ जर आणि फक्त जर $\Gamma \Entails !A$
\end{enumerate}
वरील विधानांतील ``फक्त जर'' दिशेला \emph{निर्दोषता} म्हणतात.
!!^a{derivation} पद्धत तेव्हा निर्दोष असते, जेव्हा !!{derivability} मुळे
तार्किक निष्पन्नतेची (किंवा वैधतेची) हमी मिळते. प्रत्येक समाधानकारक
!!{derivation} पद्धत निर्दोष असलीच पाहिजे; सदोष !!{derivation} पद्धती
मुळीच उपयोगी नसतात. कारण !!a{derivation} चा संपूर्ण उद्देशच वैधतेची
किंवा तार्किक निष्पन्नतेची विन्यासात्मक हमी देणे हा असतो. आपण मांडत
असलेल्या !!{derivation} पद्धतींची निर्दोषता पुढे सिद्ध करू.

उलट ``जर'' दिशेलाही महत्त्व आहे: तिला \emph{संपूर्णता} म्हणतात.
एखादी संपूर्ण !!{derivation} पद्धत इतकी समर्थ असते की $!A$~वैध असेल
तेव्हा $!A$~हे प्रमेय आहे, तसेच $\Gamma \Entails !A$ असेल तेव्हा
$\Gamma \Proves !A$, हे ती दाखवू शकते. संपूर्णता प्रस्थापित करणे अधिक
कठीण असते आणि काही तर्कप्रणालींसाठी संपूर्ण !!{derivation} पद्धतीच
नसतात. प्रथम-क्रम तर्कशास्त्रासाठी मात्र त्या आहेत. कुर्ट गोडेल यांनी
आपल्या १९२९ सालच्या प्रबंधात प्रथम-क्रम तर्कशास्त्राच्या
!!a{derivation} पद्धतीची संपूर्णता सर्वप्रथम सिद्ध केली.

!!{derivation} पद्धतींशी जोडलेली आणखी एक संकल्पना म्हणजे
\emph{सुसंगतता}. !!{sentence}s च्या एखाद्या संचापासून कोणतीही गोष्ट
!!{derive}d करता येत असेल तर तो संच विसंगत म्हणतात; अन्यथा तो सुसंगत
असतो. विसंगतता ही अपूर्ततायोग्यतेची विन्यासात्मक समकक्ष संकल्पना आहे:
जसे अपूर्ततायोग्य संच चांगल्या उपपत्ती ठरत नाहीत, तसेच !!{sentence}s चे
विसंगत संचही चांगल्या उपपत्ती ठरत नाहीत; त्यांच्यात मूलभूत स्वरूपाचा
दोष असतो. !!{sentence}s चे सुसंगत संच सत्य किंवा उपयोगी असतीलच असे
नाही, पण ते निदान तार्किक उपयोगितेचा किमान उंबरठा ओलांडतात. वेगवेगळ्या
!!{derivation} पद्धतींसाठी !!{sentence}s च्या संचांच्या सुसंगततेची नेमकी
व्याख्या वेगळी असू शकते; पण~$\Proves$ प्रमाणेच सुसंगतताही तिच्या
चिन्हार्थविषयक समकक्ष संकल्पनेशी, म्हणजे पूर्ततायोग्यतेशी, जुळली पाहिजे.
$\Gamma$ सुसंगत आहे जर आणि फक्त जर तो पूर्ततायोग्य आहे, हे नेहमी खरे
असावे अशी आपली अपेक्षा आहे. येथे ``फक्त जर'' दिशा संपूर्णतेइतकीच ठरते
(सुसंगततेमुळे पूर्ततायोग्यतेची हमी मिळते), आणि ``जर'' दिशा
निर्दोषतेइतकीच ठरते (पूर्ततायोग्यतेमुळे सुसंगततेची हमी मिळते).
प्रत्यक्षात, अभिजात प्रथम-क्रम तर्कशास्त्रासाठी निर्दोषता आणि संपूर्णता
यांची ही दोन रूपे सममूल्य आहेत.

\end{document}
